\documentclass[12pt,a4paper]{article}

\usepackage[utf8]{inputenc}
\usepackage[T1]{fontenc}
\usepackage{amsmath, amssymb}
\usepackage{enumitem}
\usepackage{geometry}
\usepackage{xcolor}
\usepackage{multicol}

\geometry{margin=2.3cm}

\title{Banco de Questões Objetivas -- Aula 12\\Processamento de Transações}
\author{Banco de Dados}
\date{}

\begin{document}

\maketitle

\section*{Instruções}

Este banco contém questões objetivas sobre os principais tópicos da Aula 12:
\begin{itemize}
    \item conceito de transação;
    \item sistemas de processamento de transações;
    \item operações de leitura e escrita;
    \item estados de uma transação;
    \item propriedades ACID;
    \item planos de execução;
    \item restaurabilidade;
    \item seriabilidade;
    \item equivalência de planos;
    \item conflitos entre operações;
    \item suporte a transações em SQL;
    \item comandos \texttt{COMMIT}, \texttt{ROLLBACK} e \texttt{START TRANSACTION}.
\end{itemize}

\section{Questões de Múltipla Escolha}

\begin{enumerate}[label=\textbf{Q\arabic*.}]

\item Uma transação em banco de dados pode ser definida como:

\begin{enumerate}[label=\alph*)]
    \item uma unidade lógica de processamento de banco de dados.
    \item uma tabela temporária sem chave primária.
    \item um comando usado apenas para criar tabelas.
    \item um índice físico do banco.
\end{enumerate}

\item Uma transação pode incluir operações de:

\begin{enumerate}[label=\alph*)]
    \item leitura, escrita, inserção, atualização e remoção.
    \item apenas criação de tabelas.
    \item apenas ordenação de resultados.
    \item apenas definição de usuários.
\end{enumerate}

\item Sistemas de processamento de transações são comuns em aplicações como:

\begin{enumerate}[label=\alph*)]
    \item bancos, reservas de passagens, cartões de crédito e compras online.
    \item apenas editores de texto locais.
    \item apenas calculadoras sem armazenamento.
    \item apenas arquivos estáticos sem usuários.
\end{enumerate}

\item Uma característica desejada em sistemas de processamento de transações é:

\begin{enumerate}[label=\alph*)]
    \item alta disponibilidade.
    \item ausência de usuários concorrentes.
    \item impossibilidade de recuperação.
    \item bloqueio permanente do banco.
\end{enumerate}

\item Outra característica desejada nesses sistemas é:

\begin{enumerate}[label=\alph*)]
    \item baixo tempo de resposta.
    \item execução exclusivamente manual.
    \item ausência de atualizações.
    \item impossibilidade de acesso simultâneo.
\end{enumerate}

\item Em um ambiente com muitos usuários, o SGBD precisa lidar com:

\begin{enumerate}[label=\alph*)]
    \item usuários concorrentes executando transações simultaneamente.
    \item apenas um usuário por semana.
    \item somente comandos de criação de esquema.
    \item ausência total de falhas.
\end{enumerate}

\item As propriedades desejáveis de transações são conhecidas pela sigla:

\begin{enumerate}[label=\alph*)]
    \item ACID.
    \item SQL.
    \item DDL.
    \item RAID.
\end{enumerate}

\item A propriedade de atomicidade significa que:

\begin{enumerate}[label=\alph*)]
    \item a transação deve ser executada completamente ou não produzir efeito algum.
    \item a transação sempre executa apenas uma operação.
    \item a transação nunca pode ser desfeita.
    \item a transação dispensa controle de concorrência.
\end{enumerate}

\item A propriedade de consistência significa que:

\begin{enumerate}[label=\alph*)]
    \item a transação deve levar o banco de um estado consistente para outro estado consistente.
    \item a transação deve ignorar todas as restrições.
    \item a transação deve gerar valores duplicados obrigatoriamente.
    \item a transação não pode alterar dados.
\end{enumerate}

\item A propriedade de isolamento significa que:

\begin{enumerate}[label=\alph*)]
    \item uma transação deve executar como se estivesse isolada das demais.
    \item uma transação nunca pode ocorrer ao mesmo tempo que outra.
    \item uma transação deve sempre visualizar dados não confirmados.
    \item uma transação deve eliminar todas as outras.
\end{enumerate}

\item A propriedade de durabilidade significa que:

\begin{enumerate}[label=\alph*)]
    \item depois de efetivada, a transação deve permanecer no banco mesmo diante de falhas.
    \item a transação deve durar muito tempo em execução.
    \item a transação deve ser apagada após o \texttt{COMMIT}.
    \item a transação só existe enquanto o usuário está conectado.
\end{enumerate}

\item O comando \texttt{COMMIT} é usado para:

\begin{enumerate}[label=\alph*)]
    \item efetivar as alterações de uma transação.
    \item desfazer as alterações de uma transação.
    \item criar uma tabela.
    \item remover uma coluna.
\end{enumerate}

\item O comando \texttt{ROLLBACK} é usado para:

\begin{enumerate}[label=\alph*)]
    \item desfazer as alterações de uma transação não efetivada.
    \item confirmar permanentemente uma transação.
    \item criar um índice.
    \item definir uma chave estrangeira.
\end{enumerate}

\item O comando \texttt{START TRANSACTION} ou \texttt{BEGIN} é usado para:

\begin{enumerate}[label=\alph*)]
    \item iniciar explicitamente uma transação.
    \item finalizar uma consulta.
    \item criar uma visão.
    \item remover uma tabela.
\end{enumerate}

\item Uma transação que terminou com sucesso normalmente chega ao estado:

\begin{enumerate}[label=\alph*)]
    \item efetivada.
    \item abortada.
    \item bloqueada permanentemente.
    \item inexistente.
\end{enumerate}

\item Uma transação que falha e tem seus efeitos desfeitos chega ao estado:

\begin{enumerate}[label=\alph*)]
    \item abortada.
    \item efetivada.
    \item parcialmente efetivada para sempre.
    \item serializável.
\end{enumerate}

\item O estado ativo de uma transação ocorre quando:

\begin{enumerate}[label=\alph*)]
    \item a transação está em execução.
    \item a transação já foi apagada do log.
    \item a transação ainda não começou e nunca começará.
    \item a transação foi substituída por uma tabela.
\end{enumerate}

\item Uma transação parcialmente efetivada é aquela que:

\begin{enumerate}[label=\alph*)]
    \item executou sua última operação, mas ainda pode precisar garantir persistência.
    \item ainda não executou nenhuma operação.
    \item foi obrigatoriamente desfeita.
    \item não possui operações de escrita.
\end{enumerate}

\item O modelo simples de transação normalmente considera operações como:

\begin{enumerate}[label=\alph*)]
    \item \texttt{read\_item} e \texttt{write\_item}.
    \item \texttt{create\_index} e \texttt{drop\_schema} apenas.
    \item \texttt{group\_by} e \texttt{order\_by} apenas.
    \item \texttt{grant} e \texttt{revoke} apenas.
\end{enumerate}

\item A operação de leitura de um item de dado consiste em:

\begin{enumerate}[label=\alph*)]
    \item copiar o valor do banco para uma variável da transação.
    \item remover o item do banco.
    \item confirmar uma transação.
    \item criar uma nova tabela.
\end{enumerate}

\item A operação de escrita de um item de dado consiste em:

\begin{enumerate}[label=\alph*)]
    \item gravar no banco um valor produzido pela transação.
    \item apenas consultar o catálogo do banco.
    \item iniciar uma conexão.
    \item ordenar o resultado de uma consulta.
\end{enumerate}

\item Um plano de execução, ou escalonamento, representa:

\begin{enumerate}[label=\alph*)]
    \item a ordem em que operações de transações são executadas.
    \item a estrutura física de uma tabela.
    \item o conjunto de domínios de atributos.
    \item o nome do banco de dados.
\end{enumerate}

\item Um plano serial é aquele em que:

\begin{enumerate}[label=\alph*)]
    \item todas as operações de cada transação aparecem consecutivamente, sem intercalação.
    \item todas as transações executam ao mesmo tempo operação por operação.
    \item toda operação é obrigatoriamente abortada.
    \item não existem comandos de leitura.
\end{enumerate}

\item Em um plano não serial:

\begin{enumerate}[label=\alph*)]
    \item operações de transações diferentes podem estar intercaladas.
    \item só existe uma transação.
    \item nenhuma transação pode escrever.
    \item todas as transações são apagadas.
\end{enumerate}

\item Ser serializável significa que:

\begin{enumerate}[label=\alph*)]
    \item o plano é equivalente a algum plano serial das mesmas transações.
    \item o plano precisa ser literalmente serial.
    \item o plano não pode ter operações de leitura.
    \item o plano sempre viola consistência.
\end{enumerate}

\item Ser serializável é diferente de ser serial porque:

\begin{enumerate}[label=\alph*)]
    \item um plano pode ter intercalação e ainda assim ser equivalente a um plano serial.
    \item todo plano serializável não possui transações.
    \item plano serializável sempre aborta.
    \item plano serial nunca é correto.
\end{enumerate}

\item Um plano serializável é considerado correto porque:

\begin{enumerate}[label=\alph*)]
    \item produz efeito equivalente a uma execução serial correta.
    \item elimina todas as transações.
    \item impede qualquer concorrência.
    \item transforma todas as operações em leituras.
\end{enumerate}

\item A principal vantagem de permitir planos serializáveis não seriais é:

\begin{enumerate}[label=\alph*)]
    \item permitir concorrência mantendo a correção.
    \item impedir acesso simultâneo.
    \item remover a necessidade de \texttt{COMMIT}.
    \item impedir uso de \texttt{ROLLBACK}.
\end{enumerate}

\item A serializabilidade é importante porque:

\begin{enumerate}[label=\alph*)]
    \item ajuda a definir quando uma execução concorrente é correta.
    \item substitui todas as restrições de integridade.
    \item elimina a necessidade de banco relacional.
    \item só serve para comandos DDL.
\end{enumerate}

\item A verificação prática da serializabilidade é difícil porque:

\begin{enumerate}[label=\alph*)]
    \item as operações intercaladas dependem do agendamento do sistema.
    \item SQL não possui comandos de consulta.
    \item não existem transações concorrentes.
    \item todos os planos são sempre seriais.
\end{enumerate}

\item Duas operações entram em conflito quando:

\begin{enumerate}[label=\alph*)]
    \item pertencem a transações diferentes, acessam o mesmo item e pelo menos uma delas é escrita.
    \item pertencem à mesma transação e acessam itens diferentes.
    \item são apenas duas leituras do mesmo item.
    \item são comandos de criação de tabela.
\end{enumerate}

\item Qual par de operações sobre o mesmo item gera conflito?

\begin{enumerate}[label=\alph*)]
    \item leitura e escrita.
    \item leitura e leitura.
    \item duas leituras.
    \item duas consultas sem item comum.
\end{enumerate}

\item Qual par de operações sobre o mesmo item também gera conflito?

\begin{enumerate}[label=\alph*)]
    \item escrita e escrita.
    \item leitura e leitura.
    \item duas operações em itens diferentes.
    \item duas operações da mesma transação apenas.
\end{enumerate}

\item Duas operações de leitura sobre o mesmo item:

\begin{enumerate}[label=\alph*)]
    \item não geram conflito.
    \item sempre geram conflito.
    \item sempre causam \texttt{ROLLBACK}.
    \item eliminam a transação.
\end{enumerate}

\item A equivalência por conflito considera:

\begin{enumerate}[label=\alph*)]
    \item a ordem relativa das operações conflitantes.
    \item apenas o nome das tabelas.
    \item apenas a quantidade de usuários conectados.
    \item apenas comandos de criação de índice.
\end{enumerate}

\item Um plano é serializável por conflito quando:

\begin{enumerate}[label=\alph*)]
    \item é equivalente por conflito a algum plano serial.
    \item não possui nenhuma operação.
    \item só possui comandos \texttt{CREATE TABLE}.
    \item todas as transações foram abortadas.
\end{enumerate}

\item A equivalência de visão é geralmente:

\begin{enumerate}[label=\alph*)]
    \item uma forma de equivalência mais geral que a equivalência por conflito.
    \item idêntica à ordenação alfabética de comandos SQL.
    \item uma técnica para criar visões SQL.
    \item uma restrição \texttt{NOT NULL}.
\end{enumerate}

\item Um plano restaurável é aquele em que:

\begin{enumerate}[label=\alph*)]
    \item uma transação só confirma depois das transações das quais leu valores também confirmarem.
    \item toda transação confirma antes de ler.
    \item transações nunca leem dados.
    \item todas as transações são sempre seriais.
\end{enumerate}

\item Um problema de um plano não restaurável é:

\begin{enumerate}[label=\alph*)]
    \item uma transação confirmar baseada em valor produzido por outra que depois aborta.
    \item ele nunca possuir concorrência.
    \item ele sempre ser serial.
    \item ele impedir qualquer leitura.
\end{enumerate}

\item Leitura suja ocorre quando:

\begin{enumerate}[label=\alph*)]
    \item uma transação lê valor escrito por outra transação ainda não confirmada.
    \item uma transação lê um valor já confirmado.
    \item uma transação apenas consulta o catálogo.
    \item uma transação executa \texttt{COMMIT}.
\end{enumerate}

\item Uma reversão em cascata pode ocorrer quando:

\begin{enumerate}[label=\alph*)]
    \item uma transação abortada força o aborto de outras que leram seus dados não confirmados.
    \item todas as transações são independentes.
    \item nenhuma transação leu dados de outra.
    \item todos os planos são seriais.
\end{enumerate}

\item Um plano sem cascata evita:

\begin{enumerate}[label=\alph*)]
    \item que transações leiam valores escritos por transações ainda não confirmadas.
    \item que transações façam leitura de dados confirmados.
    \item que transações usem \texttt{COMMIT}.
    \item que o banco tenha chaves primárias.
\end{enumerate}

\item Um plano estrito é importante porque:

\begin{enumerate}[label=\alph*)]
    \item facilita recuperação, evitando leitura ou escrita sobre itens escritos por transações não confirmadas.
    \item obriga todas as transações a serem removidas.
    \item elimina a necessidade de log em qualquer sistema.
    \item impede qualquer uso de SQL.
\end{enumerate}

\item A restaurabilidade está mais ligada à propriedade de:

\begin{enumerate}[label=\alph*)]
    \item recuperação correta após falhas ou abortos.
    \item definição de esquemas.
    \item criação de visões.
    \item normalização até BCNF.
\end{enumerate}

\item A seriabilidade está mais ligada à propriedade de:

\begin{enumerate}[label=\alph*)]
    \item isolamento.
    \item atomicidade apenas.
    \item criação de tabelas.
    \item criptografia de senha.
\end{enumerate}

\item A durabilidade está mais diretamente ligada ao uso de mecanismos como:

\begin{enumerate}[label=\alph*)]
    \item armazenamento persistente e recuperação.
    \item apenas ordenação de consultas.
    \item apenas chaves estrangeiras.
    \item apenas \texttt{SELECT DISTINCT}.
\end{enumerate}

\item A atomicidade está mais diretamente ligada ao uso de:

\begin{enumerate}[label=\alph*)]
    \item \texttt{COMMIT} e \texttt{ROLLBACK}.
    \item \texttt{GROUP BY} e \texttt{HAVING}.
    \item \texttt{CREATE VIEW} e \texttt{DROP VIEW}.
    \item \texttt{ORDER BY} e \texttt{DISTINCT}.
\end{enumerate}

\item Uma execução concorrente pode melhorar:

\begin{enumerate}[label=\alph*)]
    \item desempenho e utilização dos recursos do sistema.
    \item inconsistência obrigatória do banco.
    \item impossibilidade de recuperação.
    \item ausência de usuários simultâneos.
\end{enumerate}

\item O risco da execução concorrente sem controle é:

\begin{enumerate}[label=\alph*)]
    \item produzir resultados incorretos ou inconsistentes.
    \item tornar toda consulta serializável automaticamente.
    \item impedir qualquer conflito.
    \item eliminar necessidade de transações.
\end{enumerate}

\item Em SQL, quando o modo \texttt{autocommit} está ativado:

\begin{enumerate}[label=\alph*)]
    \item cada comando pode ser tratado como uma transação automaticamente confirmada.
    \item nenhuma transação pode ser efetivada.
    \item todo comando vira \texttt{ROLLBACK}.
    \item o banco deixa de aceitar \texttt{SELECT}.
\end{enumerate}

\item Para controlar manualmente uma transação em SQL, é comum usar:

\begin{enumerate}[label=\alph*)]
    \item \texttt{START TRANSACTION}, seguido de comandos, e depois \texttt{COMMIT} ou \texttt{ROLLBACK}.
    \item apenas \texttt{SELECT DISTINCT}.
    \item apenas \texttt{CREATE TABLE}.
    \item apenas \texttt{ORDER BY}.
\end{enumerate}

\item Em JDBC, o controle de transações pode envolver:

\begin{enumerate}[label=\alph*)]
    \item desativar \texttt{autoCommit}, executar comandos e chamar \texttt{commit} ou \texttt{rollback}.
    \item usar apenas \texttt{GROUP BY}.
    \item criar uma tabela para cada transação.
    \item remover o driver JDBC.
\end{enumerate}

\item Se uma aplicação executa duas atualizações que devem ocorrer juntas, o ideal é:

\begin{enumerate}[label=\alph*)]
    \item colocá-las na mesma transação.
    \item executar uma e ignorar a outra.
    \item usar \texttt{DROP TABLE}.
    \item armazenar tudo em uma única string.
\end{enumerate}

\item Em uma transferência bancária, debitar uma conta e creditar outra deve ser uma transação por causa da:

\begin{enumerate}[label=\alph*)]
    \item atomicidade.
    \item ausência de chaves.
    \item eliminação de \texttt{WHERE}.
    \item criação de índices.
\end{enumerate}

\item Se o banco confirma uma transação, mas perde seus dados após queda de energia, a propriedade violada é:

\begin{enumerate}[label=\alph*)]
    \item durabilidade.
    \item isolamento.
    \item 1FN.
    \item chave estrangeira.
\end{enumerate}

\item Se uma transação lê dados temporários de outra ainda não confirmada, o problema está relacionado a:

\begin{enumerate}[label=\alph*)]
    \item isolamento e leitura suja.
    \item normalização.
    \item criação de esquema.
    \item \texttt{ORDER BY}.
\end{enumerate}

\item Se uma transação deixa o banco violando uma restrição de integridade, o problema está relacionado à:

\begin{enumerate}[label=\alph*)]
    \item consistência.
    \item durabilidade apenas.
    \item criação de visão.
    \item consulta aninhada.
\end{enumerate}

\item O estudo de planos baseados em restaurabilidade busca responder:

\begin{enumerate}[label=\alph*)]
    \item se o plano permite recuperação correta diante de abortos ou falhas.
    \item se uma tabela está em BCNF.
    \item se uma consulta possui \texttt{GROUP BY}.
    \item se uma chave é composta.
\end{enumerate}

\item O estudo de planos baseados em seriabilidade busca responder:

\begin{enumerate}[label=\alph*)]
    \item se uma intercalação concorrente equivale a alguma execução serial.
    \item se uma senha está criptografada.
    \item se uma tabela possui valores nulos.
    \item se um índice é B-tree.
\end{enumerate}

\end{enumerate}

\section{Questões de Verdadeiro ou Falso}

\begin{enumerate}[resume,label=\textbf{Q\arabic*.}]

\item Julgue as afirmações:

\begin{enumerate}[label=\alph*)]
    \item ( \quad ) Uma transação é uma unidade lógica de processamento de BD.
    \item ( \quad ) Uma transação pode envolver leitura e escrita.
    \item ( \quad ) Uma aplicação pode conter várias transações.
    \item ( \quad ) Transações são usadas apenas para criar tabelas.
\end{enumerate}

\item Julgue as afirmações:

\begin{enumerate}[label=\alph*)]
    \item ( \quad ) Sistemas de transações devem lidar com usuários concorrentes.
    \item ( \quad ) Alta disponibilidade é desejável em sistemas de transações.
    \item ( \quad ) Baixo tempo de resposta é desejável em sistemas de transações.
    \item ( \quad ) Sistemas de transações não aparecem em bancos ou compras online.
\end{enumerate}

\item Julgue as afirmações:

\begin{enumerate}[label=\alph*)]
    \item ( \quad ) Atomicidade significa tudo ou nada.
    \item ( \quad ) Consistência significa preservar regras do banco.
    \item ( \quad ) Isolamento trata da interferência entre transações.
    \item ( \quad ) Durabilidade significa que a transação confirmada pode ser perdida livremente.
\end{enumerate}

\item Julgue as afirmações:

\begin{enumerate}[label=\alph*)]
    \item ( \quad ) \texttt{COMMIT} efetiva alterações.
    \item ( \quad ) \texttt{ROLLBACK} desfaz alterações não confirmadas.
    \item ( \quad ) \texttt{START TRANSACTION} pode iniciar uma transação.
    \item ( \quad ) \texttt{COMMIT} e \texttt{ROLLBACK} possuem exatamente o mesmo efeito.
\end{enumerate}

\item Julgue as afirmações:

\begin{enumerate}[label=\alph*)]
    \item ( \quad ) Um plano serial executa as operações de cada transação consecutivamente.
    \item ( \quad ) Um plano não serial pode intercalar operações de transações.
    \item ( \quad ) Um plano serializável deve ser equivalente a algum plano serial.
    \item ( \quad ) Todo plano serializável precisa ser literalmente serial.
\end{enumerate}

\item Julgue as afirmações:

\begin{enumerate}[label=\alph*)]
    \item ( \quad ) Ser serializável é diferente de ser serial.
    \item ( \quad ) Um plano serializável é considerado correto.
    \item ( \quad ) A intercalação pode ser apropriada e ainda assim gerar resultado correto.
    \item ( \quad ) A serializabilidade sempre impede concorrência.
\end{enumerate}

\item Julgue as afirmações:

\begin{enumerate}[label=\alph*)]
    \item ( \quad ) Duas leituras sobre o mesmo item não geram conflito.
    \item ( \quad ) Leitura e escrita sobre o mesmo item podem gerar conflito.
    \item ( \quad ) Escrita e escrita sobre o mesmo item podem gerar conflito.
    \item ( \quad ) Operações em itens diferentes sempre geram conflito.
\end{enumerate}

\item Julgue as afirmações:

\begin{enumerate}[label=\alph*)]
    \item ( \quad ) Equivalência por conflito preserva a ordem das operações conflitantes.
    \item ( \quad ) Um plano serializável por conflito é equivalente por conflito a algum plano serial.
    \item ( \quad ) Equivalência de visão é uma noção de equivalência entre planos.
    \item ( \quad ) Equivalência por conflito é usada para criar visões SQL.
\end{enumerate}

\item Julgue as afirmações:

\begin{enumerate}[label=\alph*)]
    \item ( \quad ) Plano restaurável evita que uma transação confirme antes daquelas das quais depende.
    \item ( \quad ) Leitura suja envolve leitura de valor não confirmado.
    \item ( \quad ) Reversão em cascata pode ocorrer por leitura de dados não confirmados.
    \item ( \quad ) Plano sem cascata permite leitura de dados não confirmados livremente.
\end{enumerate}

\item Julgue as afirmações:

\begin{enumerate}[label=\alph*)]
    \item ( \quad ) Plano estrito facilita recuperação.
    \item ( \quad ) Restaurabilidade está relacionada a recuperação correta.
    \item ( \quad ) Seriabilidade está relacionada à correção da execução concorrente.
    \item ( \quad ) Planos não restauráveis são sempre desejáveis.
\end{enumerate}

\item Julgue as afirmações:

\begin{enumerate}[label=\alph*)]
    \item ( \quad ) Concorrência pode melhorar desempenho.
    \item ( \quad ) Concorrência sem controle pode causar inconsistências.
    \item ( \quad ) Controle de transações ajuda a preservar propriedades ACID.
    \item ( \quad ) Concorrência elimina a necessidade de isolamento.
\end{enumerate}

\item Julgue as afirmações:

\begin{enumerate}[label=\alph*)]
    \item ( \quad ) No modo \texttt{autocommit}, comandos podem ser confirmados automaticamente.
    \item ( \quad ) Em SQL, \texttt{BEGIN} pode iniciar uma transação.
    \item ( \quad ) Em JDBC, pode-se usar \texttt{commit} e \texttt{rollback}.
    \item ( \quad ) Transações não podem ser controladas por aplicações.
\end{enumerate}

\end{enumerate}

\section{Questões de Aplicação Objetiva}

\begin{enumerate}[resume,label=\textbf{Q\arabic*.}]

\item Considere uma transferência bancária com duas operações: debitar R\$100 da conta A e creditar R\$100 na conta B. Se o débito ocorre, mas o crédito não ocorre, qual propriedade foi violada?

\begin{enumerate}[label=\alph*)]
    \item Atomicidade.
    \item Durabilidade.
    \item Isolamento apenas.
    \item BCNF.
\end{enumerate}

\item Uma transação confirmou a compra de uma passagem. Logo depois, o sistema caiu e a compra desapareceu. Qual propriedade foi violada?

\begin{enumerate}[label=\alph*)]
    \item Durabilidade.
    \item Atomicidade.
    \item Isolamento.
    \item 2FN.
\end{enumerate}

\item Duas transações executam ao mesmo tempo sobre o mesmo saldo bancário e uma interfere na outra, produzindo resultado incorreto. O problema está mais diretamente ligado à falta de:

\begin{enumerate}[label=\alph*)]
    \item isolamento.
    \item chave estrangeira.
    \item normalização.
    \item visão SQL.
\end{enumerate}

\item Uma transação finaliza deixando o banco violando uma regra de integridade. A propriedade afetada é:

\begin{enumerate}[label=\alph*)]
    \item consistência.
    \item durabilidade.
    \item ordenação.
    \item projeção.
\end{enumerate}

\item Considere o plano:

\[
r_1(X),\ w_1(X),\ r_2(X),\ w_2(X)
\]

As operações $w_1(X)$ e $r_2(X)$:

\begin{enumerate}[label=\alph*)]
    \item entram em conflito, pois acessam o mesmo item e uma é escrita.
    \item não entram em conflito, pois são de transações diferentes.
    \item não entram em conflito, pois leitura e escrita nunca conflitam.
    \item são comandos DDL.
\end{enumerate}

\item Considere o plano:

\[
r_1(X),\ r_2(X)
\]

Essas operações:

\begin{enumerate}[label=\alph*)]
    \item não entram em conflito.
    \item entram em conflito sempre.
    \item causam \texttt{ROLLBACK} obrigatório.
    \item violam durabilidade.
\end{enumerate}

\item Considere duas transações $T_1$ e $T_2$. Um plano executa todas as operações de $T_1$ e depois todas as operações de $T_2$. Esse plano é:

\begin{enumerate}[label=\alph*)]
    \item serial.
    \item não restaurável obrigatoriamente.
    \item não serializável obrigatoriamente.
    \item uma leitura suja.
\end{enumerate}

\item Um plano intercala operações de $T_1$ e $T_2$, mas produz o mesmo efeito de executar $T_1$ antes de $T_2$. Esse plano pode ser:

\begin{enumerate}[label=\alph*)]
    \item serializável.
    \item necessariamente incorreto.
    \item necessariamente não restaurável.
    \item uma operação DDL.
\end{enumerate}

\item $T_2$ lê um valor escrito por $T_1$. Depois $T_2$ confirma antes de $T_1$. Se $T_1$ abortar, o plano é problemático porque:

\begin{enumerate}[label=\alph*)]
    \item $T_2$ confirmou baseada em dado não confirmado.
    \item $T_2$ nunca leu dados.
    \item $T_1$ e $T_2$ eram a mesma transação.
    \item não houve nenhuma dependência.
\end{enumerate}

\item Qual sequência SQL representa melhor uma transação manual?

\begin{enumerate}[label=\alph*)]
    \item \texttt{START TRANSACTION;} comandos SQL; \texttt{COMMIT;}
    \item \texttt{GROUP BY;} comandos SQL; \texttt{ORDER BY;}
    \item \texttt{CREATE TABLE;} \texttt{DROP TABLE;} \texttt{SELECT;}
    \item \texttt{WHERE;} \texttt{HAVING;} \texttt{DISTINCT;}
\end{enumerate}

\item Considere o plano:
\[
r_1(X),\ w_1(X),\ r_2(X),\ w_2(X)
\]
Qual conflito ocorre entre $w_1(X)$ e $r_2(X)$?

\begin{enumerate}[label=\alph*)]
    \item Não há conflito, pois são transações diferentes.
    \item Conflito do tipo escrita-leitura no mesmo item.
    \item Conflito apenas se $X$ for chave primária.
    \item Conflito de definição de tabela.
\end{enumerate}

\item Considere o plano:
\[
r_1(X),\ r_2(X),\ w_1(X),\ w_2(X)
\]
Quais pares de operações sobre $X$ entram em conflito?

\begin{enumerate}[label=\alph*)]
    \item Apenas $r_1(X)$ e $r_2(X)$.
    \item $r_1(X)$ com $w_2(X)$, $r_2(X)$ com $w_1(X)$ e $w_1(X)$ com $w_2(X)$.
    \item Nenhum par, pois todas usam o mesmo item.
    \item Apenas leituras entram em conflito.
\end{enumerate}

\item Considere o plano:
\[
r_1(X),\ w_1(X),\ r_2(X),\ w_2(X)
\]
No grafo de precedência, qual aresta deve existir?

\begin{enumerate}[label=\alph*)]
    \item $T_2 \rightarrow T_1$ apenas.
    \item $T_1 \rightarrow T_2$.
    \item Nenhuma, pois o plano é serial.
    \item $X \rightarrow T_1$.
\end{enumerate}

\item Considere o plano:
\[
r_1(X),\ w_2(X),\ w_1(X)
\]
Qual relação de precedência é gerada pelos conflitos?

\begin{enumerate}[label=\alph*)]
    \item Apenas $T_1 \rightarrow T_2$.
    \item Apenas $T_2 \rightarrow T_1$.
    \item Há arestas nos dois sentidos, indicando ciclo.
    \item Não há conflitos.
\end{enumerate}

\item Um plano cujo grafo de precedência possui ciclo é:

\begin{enumerate}[label=\alph*)]
    \item serializável por conflito.
    \item não serializável por conflito.
    \item obrigatoriamente serial.
    \item sempre sem cascata.
\end{enumerate}

\item Considere o plano:
\[
r_1(X),\ w_1(X),\ r_2(Y),\ w_2(Y)
\]
Assumindo que $X$ e $Y$ são itens distintos, esse plano é:

\begin{enumerate}[label=\alph*)]
    \item não serializável por conflito, pois é intercalado.
    \item serializável por conflito, pois não há conflitos entre $T_1$ e $T_2$.
    \item obrigatoriamente não restaurável.
    \item um exemplo de leitura suja.
\end{enumerate}

\item Considere o plano:
\[
w_1(X),\ r_2(X),\ c_2,\ c_1
\]
Se $c_i$ representa \texttt{COMMIT} de $T_i$, esse plano é:

\begin{enumerate}[label=\alph*)]
    \item restaurável, pois $T_2$ confirma antes de $T_1$.
    \item não restaurável, pois $T_2$ leu valor escrito por $T_1$ e confirmou antes de $T_1$.
    \item estrito, pois todo plano com \texttt{COMMIT} é estrito.
    \item serial, pois possui duas confirmações.
\end{enumerate}

\item Considere o plano:
\[
w_1(X),\ r_2(X),\ c_1,\ c_2
\]
Esse plano é restaurável porque:

\begin{enumerate}[label=\alph*)]
    \item $T_2$ confirma depois da transação de quem leu o valor.
    \item $T_2$ confirma antes de $T_1$.
    \item não existe leitura de $X$.
    \item toda escrita gera rollback automático.
\end{enumerate}

\item Considere o plano:
\[
w_1(X),\ r_2(X),\ c_1,\ c_2
\]
Esse plano é sem cascata?

\begin{enumerate}[label=\alph*)]
    \item Sim, pois $T_2$ só leu $X$ depois do \texttt{COMMIT} de $T_1$.
    \item Não, pois $T_2$ leu $X$ antes do \texttt{COMMIT} de $T_1$.
    \item Sim, pois todo plano restaurável é sem cascata.
    \item Não, pois não possui \texttt{ROLLBACK}.
\end{enumerate}

\item Considere o plano:
\[
w_1(X),\ c_1,\ r_2(X),\ c_2
\]
Esse plano é:

\begin{enumerate}[label=\alph*)]
    \item sem cascata, pois $T_2$ lê o valor de $T_1$ apenas após $T_1$ confirmar.
    \item não restaurável, pois $T_2$ lê depois do \texttt{COMMIT}.
    \item sempre não serializável.
    \item inválido por conter leitura.
\end{enumerate}

\item Considere o plano:
\[
w_1(X),\ w_2(X),\ c_1,\ c_2
\]
Esse plano viola a propriedade de plano estrito porque:

\begin{enumerate}[label=\alph*)]
    \item $T_2$ escreve $X$ antes de $T_1$ confirmar ou abortar.
    \item duas escritas nunca entram em conflito.
    \item $T_1$ confirma antes de $T_2$.
    \item não existe operação de leitura.
\end{enumerate}

\item Considere o plano:
\[
w_1(X),\ a_1,\ r_2(X)
\]
Se $a_1$ representa \texttt{ABORT} de $T_1$, a leitura de $T_2$:

\begin{enumerate}[label=\alph*)]
    \item ocorre depois do abort de $T_1$, então não lê valor não confirmado de $T_1$.
    \item é sempre leitura suja.
    \item torna o plano não serializável automaticamente.
    \item equivale a \texttt{COMMIT}.
\end{enumerate}

\item Considere a situação: duas transações leem o mesmo saldo $X=100$, ambas calculam novo valor e ambas escrevem no banco, uma com $X=90$ e outra com $X=80$. O problema típico é:

\begin{enumerate}[label=\alph*)]
    \item atualização perdida.
    \item normalização excessiva.
    \item chave estrangeira inválida.
    \item cobertura mínima.
\end{enumerate}

\item Considere:
\[
r_1(X),\ w_1(X),\ r_2(X),\ a_1
\]
Se $T_2$ leu o valor escrito por $T_1$ antes de $T_1$ abortar, ocorreu:

\begin{enumerate}[label=\alph*)]
    \item leitura suja.
    \item leitura serial.
    \item plano estrito.
    \item 1FN.
\end{enumerate}

\item Qual sequência representa corretamente uma transação SQL que desfaz alterações ao detectar erro?

\begin{enumerate}[label=\alph*)]
    \item \texttt{START TRANSACTION;} comandos; erro; \texttt{ROLLBACK;}
    \item \texttt{COMMIT;} comandos; erro; \texttt{START TRANSACTION;}
    \item \texttt{GROUP BY;} comandos; erro; \texttt{HAVING;}
    \item \texttt{CREATE TABLE;} erro; \texttt{SELECT DISTINCT;}
\end{enumerate}

\item Em um plano restaurável, se $T_j$ lê um item escrito por $T_i$, então:

\begin{enumerate}[label=\alph*)]
    \item $T_j$ deve confirmar antes de $T_i$.
    \item $T_j$ só deve confirmar depois de $T_i$ confirmar.
    \item $T_i$ nunca pode escrever.
    \item $T_j$ deve obrigatoriamente abortar.
\end{enumerate}

\item Em um plano sem cascata, uma transação só pode ler valores escritos por outra transação depois que:

\begin{enumerate}[label=\alph*)]
    \item a outra transação confirmar.
    \item a outra transação iniciar.
    \item a outra transação ler qualquer item.
    \item o banco ser normalizado.
\end{enumerate}

\item Em um plano estrito, uma transação não deve ler nem escrever um item $X$ escrito por outra transação até que essa outra transação:

\begin{enumerate}[label=\alph*)]
    \item confirme ou aborte.
    \item leia $X$ novamente.
    \item crie uma tabela.
    \item execute \texttt{SELECT DISTINCT}.
\end{enumerate}


\end{enumerate}

\newpage

\section*{Gabarito}

\begin{multicols}{2}
\begin{enumerate}[label=\textbf{Q\arabic*.}]
    \item a
    \item a
    \item a
    \item a
    \item a
    \item a
    \item a
    \item a
    \item a
    \item a
    \item a
    \item a
    \item a
    \item a
    \item a
    \item a
    \item a
    \item a
    \item a
    \item a
    \item a
    \item a
    \item a
    \item a
    \item a
    \item a
    \item a
    \item a
    \item a
    \item a
    \item a
    \item a
    \item a
    \item a
    \item a
    \item a
    \item a
    \item a
    \item a
    \item a
    \item a
    \item a
    \item a
    \item a
    \item a
    \item a
    \item a
    \item a
    \item a
    \item a
    \item a
    \item a
    \item a
    \item a
    \item a
    \item a
    \item a
    \item a
    \item a
    \item a
    \item a
    \item a
    \item a
    \item a
    \item a
    \item a
    \item a
    \item V, V, V, F
    \item V, V, V, F
    \item V, V, V, F
    \item V, V, V, F
    \item V, V, V, F
    \item V, V, V, F
    \item V, V, V, F
    \item V, V, V, F
    \item V, V, V, F
    \item V, V, V, F
    \item V, V, V, F
    \item V, V, V, F
    \item a
    \item a
    \item a
    \item a
    \item a
    \item a
    \item a
    \item a
    \item a
    \item a
    \item b
    \item b
    \item b
    \item c
    \item b
    \item b
    \item b
    \item a
    \item b
    \item a
    \item a
    \item a
    \item a
    \item a
    \item a
    \item b
    \item a
    \item a
\end{enumerate}
\end{multicols}

\end{document}
